\section{Речник на използваните съкращения и понятия}
\label{sec:glossary}

Настоящият раздел има две функции. От една страна, той обобщава използваните съкращения и термини, за да улесни последователното четене на текста. От друга, той фиксира точния смисъл, в който някои понятия се употребяват в рамките на дипломната работа. Това е важно, защото думи като ``модел'', ``конектор'', ``migration'' или ``capability'' могат да имат различни значения в зависимост от контекста.

\subsection{Основни съкращения}

Таблица~\ref{tab:glossary} обобщава основните съкращения и термини, използвани в дипломната работа.

\begin{table}[H]
\centering
\caption{Речник на използваните съкращения}
\label{tab:glossary}
\begin{tabular}{>{\bfseries}l L{5.8cm} L{6.0cm}}
\toprule
\textbf{Съкращение} & \textbf{Пълна форма} & \textbf{Контекст} \\
\midrule
API & Application Programming Interface & Програмен интерфейс между софтуерни компоненти \\
BSON & Binary JSON & Бинарен формат за сериализация, използван от MongoDB \\
CQL & Cassandra Query Language & Език за заявки към Apache Cassandra \\
CRUD & Create, Read, Update, Delete & Четирите основни операции върху данни \\
DDL & Data Definition Language & Език за описание и промяна на схема \\
IR & Intermediate Representation & Междинно представяне на query структурата \\
ORM & Object-Relational Mapping & Картографиране между обектен и персистентен модел \\
PFR & Plain Fields Reflection & \code{Boost.PFR} механизъм за работа с агрегати \\
RAII & Resource Acquisition Is Initialization & Идиом за управление на ресурси в C++ \\
RDBMS & Relational Database Management System & Релационна система за управление на бази данни \\
SFINAE & Substitution Failure Is Not An Error & Класически механизъм за шаблонни ограничения в C++ \\
SQL & Structured Query Language & Стандартизиран език за релационни заявки \\
TMP & Template Metaprogramming & Мета-програмиране чрез шаблони \\
WG21 & Working Group 21 & Работна група към ISO за стандарта на C++ \\
\bottomrule
\end{tabular}
\end{table}

\subsection{Работни дефиниции на ключови понятия}

В рамките на дипломната работа няколко термина имат по-специфичен смисъл от общоупотребимия. Таблица~\ref{tab:working_definitions} фиксира този смисъл, за да не се смесват архитектурните и общите значения на понятията.

\begin{table}[H]
\centering
\small
\caption{Работни дефиниции на ключовите понятия в дипломната работа}
\label{tab:working_definitions}
\begin{tabularx}{\textwidth}{L{3.4cm}Y}
\toprule
\textbf{Понятие} & \textbf{Работна дефиниция} \\
\midrule
Логически модел & compile-time описанието на entity типовете, техните свойства и връзки, независимо от физическия backend \\
Query IR & типизираното междинно представяне на заявката преди конкретната материализация към SQL, BSON-подобен филтър, Redis команда, Cypher или CQL \\
Connector contract & формалният набор от изисквания към \code{connector\_trait<DB>}, чрез който backend-ът участва в системата \\
Capability tag & compile-time маркер, който ограничава допустимите операции и изразява архитектурни възможности на backend-а \\
Prepared query & обект, който свързва предварително конструиран query IR с конкретна \code{db<DB>} инстанция за многократно изпълнение \\
Schema-aware поддръжка & ниво на интеграция, при което библиотеката може да работи със schema metadata, DDL hooks и migration diff логика \\
Hydration & материализиране на резултатите от backend-а в типизиран C++ резултатен обект \\
Verification & съвместната роля на compile-time ограниченията, unit тестовете и integration тестовете за доказване на коректност \\
\bottomrule
\end{tabularx}
\end{table}

\subsection{Нотация и именуване в текста}

Дипломната работа използва и последователна нотация за означаване на типове, traits и backend роли. Тази нотация не е случайна; тя отразява реалната структура на хранилището и улеснява връзката между текста и кода.

\begin{table}[H]
\centering
\small
\caption{Нотация и именуване, използвани в текста}
\label{tab:naming_conventions_glossary}
\begin{tabularx}{\textwidth}{L{3.6cm}Y}
\toprule
\textbf{Означение} & \textbf{Смисъл в дипломната работа} \\
\midrule
\code{DB} & generic backend тип, използван като параметър в traits, helper-и и \code{db<DB>} обвивката \\
\code{db<DB>} & фасаден тип, който предоставя user-facing интерфейса за изпълнение на заявки \\
\code{connector\_trait} & специализацията, която описва как конкретният backend изпълнява query IR-а \\
\code{wire\_type<T>} & mapping от C++ тип към wire-level представяне за конкретния backend \\
\code{cursor\_type} & абстракция за lifecycle-а или обхождането на резултатите \\
\code{field<...>} & compile-time обозначение на избрано property поле в query API-то \\
\code{Placeholder<T>} & анонимен runtime параметър в query IR-а \\
\code{*LiveDB} & naming convention за конектори, които използват реален driver-backed backend, а не mock вариант \\
\bottomrule
\end{tabularx}
\end{table}

\subsection{Репозиториен контекст на най-често цитираните артефакти}

Тъй като дипломната работа е силно обвързана с реалното хранилище, полезно е да се фиксира и ролята на най-често цитираните файлови групи. Таблица~\ref{tab:repository_artifact_map} служи като кратък ориентир за това къде в проекта се намират основните доказателствени източници.

\begin{table}[H]
\centering
\small
\caption{Карта на основните репозиторийни артефакти, използвани в текста}
\label{tab:repository_artifact_map}
\begin{tabularx}{\textwidth}{L{3.8cm}L{3.1cm}Y}
\toprule
\textbf{Път или група файлове} & \textbf{Роля} & \textbf{Как се използва в дипломната работа} \\
\midrule
\path{lib/include/ORM/entity/} & entity модел & използва се като източник за \code{property}, \code{relationship} и логическото описание на данните \\
\path{lib/include/ORM/query/} & query IR и fluent API & използва се за анализ на \code{select}, \code{insert}, \code{update}, \code{deleteq} и query composition слоя \\
\path{lib/include/ORM/connector/} & connector contract и execution layer & използва се за формализиране на capability механизма, prepared queries и \code{db<DB>} фасадата \\
\path{lib/include/ORM/db/migration/} & migration и schema-aware логика & използва се като база за анализа на DDL hooks и migration diff pipeline-а \\
\path{tests/unit/} & unit verification & служи като доказателствен слой за локалната коректност на отделните подсистеми \\
\path{tests/integration/} & integration verification & служи за демонстрация, че абстракцията работи при реално или mock изпълнение на заявки \\
\path{doc/v2/bg/} и \path{doc/v2/en/} & документационен слой & използват се за синхронно академично представяне на архитектурата и резултатите \\
\bottomrule
\end{tabularx}
\end{table}

\subsection{Извод}

Разширеният речник показва, че терминологичната последователност е неразделна част от качеството на настоящата дипломна работа. Когато архитектурата е силно типизирана и multi-backend насочена, прецизната употреба на термините не е само стилов въпрос, а условие за коректно разбиране на цялата аргументация.
